%!TEX root = main.tex

 
We fix a fresh symbol $\separator\notin\Sigma$ and let $\Sigma':=\Sigma \cup \{\separator\}$.
Each string sequence  $s = (u_1, \cdots, u_m)$ over $\Sigma$ is encoded as a string  $\enc(s):=\separator u_1\separator \cdots \separator u_m \separator$ over $\Sigma'$. % the (extended) alphabet $\Sigma \cup \{\separator\}$.  
Note that the strings %encodings of string sequences are not arbitrary strings from $(\Sigma \cup \{\separator\})^*$, but the strings that match 
resulted from the encoding should match the regular expression $\separator (\Sigma^* \separator)^*$ and the string $\separator$ encodes the empty sequence.
%For a sequence $s$, let $\enc(s)$ denote the string encoding of $s$. 

With this encoding of sequences as strings, each string sequence operation is also transformed into the corresponding \emph{string operation}. 
 %all the sequence operations can be turned into string operations. 
%
\begin{description}
\item[sequence concatenation $s_1 \concat s_2$:] %is turned into the 
string concatenation  $\enc(s_1) \concat \enc(s_2)$. 
%
\item[sequence  write {$s[i \rightarrow u]$}: ] 
%string operation 
$\mywrite(\enc(s), i+1, u)$, which replaces the substring of $\enc(s)$ between the $(i+1)$-th occurrence of $\separator$ and the $(i+2)$-th occurrence of $\separator$, by $u$, if $i\in[0,n-2]$, where $n$ is the number of occurrences of $\separator$ in $\enc(s)$, and is undefined otherwise (i.e. $i\not\in[0,n-2]$). Note that this operation is not a typical string replace operation.  (Recall that the indices of elements in sequences start from $0$ and $n$, i.e. the number of occurrences of $\separator$ in $\enc(s)$, is one more than the length of $s$.) 
%Although this operation is not a typical replace operation, we can show that its pre-image can still be effectively computed. 

\item[sequence filter operation $\myfilter_e(s)$:] %string operation 
$\myfilter_e(\enc(s))$, which removes every substring of $\enc(s)$ between two consecutive occurrences of $\separator$ that does not match the regular expression $e$. (For instance, if $e = a^* b$ and $s = (ab, ac, aab)$, then $\myfilter_e(\enc(s)) = \separator ab \separator aab \separator$.)
%
\item[subsequence operation {$s[i, j]$}:] %$s[i, j]$ is turned into 
$\subseq(\enc(s), i+1, j)$, which keeps only the substring of $\enc(s)$ between the $(i+1)$-th occurrence and the $\min(i+j, n-1)$-th occurrence of $\separator$, if $i\in[0,n-2]$, where $n$ is the number of occurrences of $\separator$ in $\enc(s)$, and  is undefined otherwise (i.e. $i\not\in[0,n-2]$). 
%
\item [$\mysplit_e(u)$:] string operation $\mysplitstr_e(u)$ that transforms $u = u'_1 v_1 u'_2 \cdots u'_{k-1} v_{k-1} u'_k$ into  $\separator u'_1 \separator u'_2 \cdots u'_{k-1} \separator u'_k \separator$, where for each $i \in [k-1]$,  $v_1$ is the leftmost and longest matching of $e$ in $u'_{i} v_i u'_{i+1}$, moreover,  $u'_k$ does not contain any substring that matches $e$. 
%
%(Note that here we abuse the notation a bit by using $\mysplit_e$ to denote both the operation of splitting a string into a sequence and the operation of transducing a string into another string. Similarly for the operation $\mymatchall_e$ in the sequel.)
%
\item[$\mymatchall_e(u)$:] %string operation 
$\mymatchallstr_e(u)$ that transforms $u = u'_1 v_1 u'_2 v_2 \cdots u'_{k-1}v_{k-1} u'_k$ into 
$\separator v_1 \separator \cdots \separator v_{k-1} \separator$, where for each $i \in [k-1]$, $v_i$ is the leftmost and longest occurrence of $e$ in $u'_i v_i u'_{i+1}$, moreover, $u'_k$ does not contain any substring that matches $e$.

%\item The transduction operation $\transducer(\seqvarb)$ is turned into a string operation $\transducer_{\separator}(\enc(\seqvarb))$, where $\transducer_{\separator}$ is the transducer that replaces every substring between  two occurrences of $\separator$, say $u$, by $\transducer(u)$.
%
\item[sequence read operation $\myat(s, i+1)$:] $\elem(\enc(s), i+1)$ that keeps only the substring between the $(i+1)$-th occurrence  and the $(i+2)$-th occurrence of $\separator$, if $i\in[0,n-2]$, where $n$ is the number of occurrences of $\separator$ in $\enc(s)$, and is undefined otherwise (i.e. $i\not\in[0,n-2]$). (For instance, if $s = (ab, ac)$ and $i = 0$, then $\elem(\enc(s), 1) = \elem(\separator ab \separator ac \separator, 1) = ab$.)
%
\item[sequence join operation $\myjoin_u(s)$:] $\myjoin_u(\enc(s))$ that removes the first and the last occurrences of $\separator$ and replaces all the remaining occurrences of $\separator$ by $u$.
\item[sequence length operation $\myseqlen(s)$:] $\myseqlen(\enc(s))$ - 1 where $\myseqlen(\enc(s))$ counts the number of occurrences of $\separator$ in $\enc(s)$.
\end{description}

%After encoding string sequences into strings, 
Hence we obtain atomic string constraints %where the string equalities are in one 
of the following forms, 
\begin{equation*}\label{eqn-str-enc}
\begin{array} {l}
\svarx = u \mid \svarx = \svary \mid \svarx = \svary \concat \svarz \mid \svarx = \transducer(\svary) \mid \svarz = \mywrite(\svarx, it, \svary) \mid  \svary  = \myfilter_e(\svarx)  \mid   \\
\svary = \mysplitstr_e(\svarx) \mid \svary = \subseq(\svarx, it_1, it_2) \mid   
 \svary = \mymatchallstr_e(\svarx) \mid \svary = \elem(\svarx, it)  \mid \\
 \svary = \myjoin_u(\svarx) \mid it = \mystrlen(\svarx) \mid it = \myseqlen(\svarx),
\end{array}
\end{equation*}
where $\transducer$ is a finite-state transducer. 

We denote by $\exstrlogic$ %(where $\sf Ex$ means ``Extended'') to denote 
the class of string constraints which are %defined as 
positive Boolean combinations (no negation) of the atomic string constraints, and  %aforementioned string equalities. Moreover, let us use 
$\slexstrlogic$ denotes its straight-line fragment.  % of $\exstrlogic$, 
%where string inequalities do not occur and there are no cyclic dependencies between string variables. 
For each formula $\varphi$ in $\slarraylogic$, $\enc(\varphi)$ is the resulting formula in $\exstrlogic$. We have the following result.

\begin{proposition} %\tl{this is to be improved}
	% $\arraylogic$ constraints into $\exstrlogic$ preserves the satisfiability. 
For each  constraint $\varphi$ in $\arraylogic$, 
	$\varphi$ and $\enc(\varphi)$ are equisatisfiable.
	Moreover, if $\varphi$ is in $\slarraylogic$, then $\enc(\varphi)$ is in $\slexstrlogic$.
\end{proposition}
%It is not hard to see that the aforementioned encoding of 
\begin{example}
Consider the $\slarraylogic$ constraint $\varphi:= (\svars_1 = \svars_0\cdot(u,v)\wedge\myseqlen(\svars_1) < 2)$ where $\svars_0$ is a string variable and $u,v$ are string constants. 
Then, its string encoding is $\enc(\varphi) := (\enc(\svars_1) = \enc(\svars_0)\cdot \separator u \separator v \separator \wedge \myseqlen(\enc(\svars_1)) - 1  < 2)$, which is unsatisfiable, since there are at least three occurrences of $\separator$ in $\enc(\svars_1)$. \qed
\end{example}

The satisfiability of $\slarraylogic$ is now reduced to that of $\slexstrlogic$. %Our goal is then to apply the framework in~\cite{atva2020} to solve the $\slexstrlogic$ constraints. At first, 
It is easy to observe that the operations $\myfilter_e$, $\mysplitstr_e$, $\mymatchallstr_e$ and $\myjoin_u$ can all be represented by a finite-state transducer, %seen as special cases of transducers $\transducer$, thus their pre-images can be effectively computed according to the 
so the results in~\cite{atva2020} are directly applicable. 
In the next section, %It remains to 
we show that the pre-images of the remaining string operations, i.e., %that are not captured directly by the operations in~\cite{atva2020}, that is, 
$\mywrite$, $\subseq$, $\elem$ and $\myseqlen$, 
%satisfy that their  
can be effectively computed. 


